\ifdefined\XeTeXversion\else\pdfoutput=1\fi
\documentclass[10pt]{nvidiaextabstract}
\usepackage{xurl}
\ifdefined\XeTeXversion\microtypesetup{tracking=false}\fi
\venue{From Human Videos to Dexterous Robot Capability (CoRL 2026)}
\submissionID{longbot -- Track 2}
\title{Evidence-Constrained Reconstruction Repair\\and Residual Humanoid Grounding}
\shorttitle{longbot: V2D Track 2}
\author[1,2]{Ismail Syed}
\affil[1]{University of Texas at Austin}
\affil[2]{\href{https://www.longhorizon.bot}{www.longhorizon.bot}}
\begin{abstract}
We present our V2D Track 2 system across clean, corrupted and RGB-reconstructed
human--object references. A common G1 pipeline combines retargeting, frozen
SONIC body control, phase/proximity hand closing and optional episode-specific
FlashSAC residual actors. Repairs are admitted using source evidence, object
settling and unassisted physical evaluation. Selected 30-episode submissions
achieve AUC 0.32917, 0.25073 and 0.19051 in Tiers 1--3, respectively. Both
success rates reach 4/30 in Tier 1 and 1/30 in Tier 2, but remain zero in Tier 3.
Our contribution is a reusable, auditable repair-to-control workflow and an
empirical account of where better visual tracking fails to produce manipulation.
\end{abstract}
\begin{document}
\maketitle
\section{Introduction}
Human demonstrations offer inexpensive supervision, but geometric uncertainty,
morphology differences and physical contact complicate transfer to a robot.
V2D~\cite{v2d} separates these effects using clean references (Tier 1), synthetically
corrupted references (Tier 2) and participant reconstructions (Tier 3).
All tiers are evaluated against clean object trajectories. We keep tier-specific
references and learned actors separate; we do not transfer clean-trained
episode actors into the noisy tiers.

We build on SONIC~\cite{sonic} and FlashCHORD/FlashSAC~\cite{flashchord,flashsac}.
Our work develops reconstruction repair, scene admission and controlled policy
selection rather than a new foundation model. The selected files are Tier 1
Submission 7, Tier 2 Submission 2 and Tier 3 Submission 4.

\section{Approach}
\textbf{Reconstruct and repair.}
Tier 3 uses a CARI4D-style initialization~\cite{cari4d}: SAM2 masks,
MoGe-2 depth, SAM 3D body/object estimates and FoundationPose object tracking.
We repair incorrect object identity, incomplete masks, metric inconsistencies,
occlusion recovery, rotation jumps and support geometry. Visible-frame anchors
and bounded temporal fitting preserve accepted intervals and test joins.
Image overlap alone cannot establish depth, contact or physical validity.
Hidden intervals retain explicit uncertainty. The permitted first clean frame
is used for alignment; later clean trajectories are reserved for evaluation,
not reconstruction or actor observations. Episode-specific proposals and manual
visual review remain part of this system.

\textbf{Retarget and control.}
MHR-to-SOMA conversion and constrained inverse kinematics generate G1 reference
motion. A portable scene includes reconstructed geometry, support, inertia and
collision assets. Frozen SONIC controls the body. Each hand closes gradually
during its reference contact phase, gated by actual palm proximity to a target
attached to the current object pose. Targets obey rate and joint limits.
Optional residual actors add bounded body and finger corrections; zero residual
does not mean zero finger targets. Each sequence has its own policy assignment.

\textbf{Train and admit.}
FlashSAC learns off-policy from bounded replay in parallel simulation. Rewards
combine body, palm and object tracking, contact guidance and regularization;
orientation and broader approach shaping are explicit experiment variants.
Curricula mix starts at frame zero with starts partway through the reference.
Early runs removed artificial object assistance around 6M transitions; recent
pilots taper it to zero by 10M. All reported evaluations start at frame zero
without assistance. Source checks, isolated object settling and repeated
rollouts precede candidate admission. A checkpoint is not selected solely
because training reward or training duration increases.

\section{Experiments \& Results}
We used a 24 GB RTX 5090 laptop and bounded L4 cloud jobs. Local GPU clocks were
capped at 1500 MHz; L4 power at 60 W. Recent jobs pause at 82 degrees C, resume at
74 degrees C and stop at 83 degrees C. Table~\ref{tab:results} reports exact Kaggle
scores for complete files, not cherry-picked repeats. AUC and success rates
are higher-is-better; MPPE is lower-is-better in centimetres.
\begin{table}[t]
\centering\small\setlength{\tabcolsep}{3pt}
\begin{tabular}{llrrrr}
\toprule
Tier & File & AUC & SP-SR & MP-SR & MPPE\\
\midrule
1 & Initial & .21646 & .00000 & .00000 & 56.78341\\
1 & Selected S7 & .32917 & .13333 & .13333 & 39.21800\\
2 & Initial & .22281 & .00000 & .00000 & 59.14549\\
2 & Selected S2 & .25073 & .03333 & .03333 & 55.10178\\
3 & Initial & .13370 & .00000 & .00000 & 70.87785\\
3 & Selected S4 & .19051 & .00000 & .00000 & 59.96399\\
\bottomrule
\end{tabular}
\caption{Each file covers all 30 episodes. Tier differences include both input
quality and policy selection; this is not a controlled noise ablation.}
\label{tab:results}
\end{table}

Tier 1 gains success from residual learning and broader episode selection.
Tier 2's selected Episode 0 actor passes both success criteria in five
confirmation rollouts. Tier 3 S4 replaces only Episodes 20 and 26 with repaired
frozen policies, retaining the other 28 designated trajectories on the same
evaluator build. Its selected set includes six previously learned residual
actors, but no episode passes the task-success criteria.

Negative results constrain interpretation. Episode 26's new 15M approach actor
extends upright survival from 9.79 to 18.51 seconds and improves reaching, yet
raises clean-target MPPE from approximately 7.96 to 9.45 cm. We retain the
repaired frozen controller. New Episode 20/7 actors can also lose to retained
policies despite better reference-relative diagnostics. Quieter object rest
can improve AUC without grasping; we do not describe these gains as successful
manipulation. Repeated rollouts measure execution variability, not variation
across independently trained seeds.

\section{Conclusion}
Conservative repairs and selective residual learning improve object tracking,
while sustained contact and task timing remain the main limitations. Our
\textbf{economic viability} contribution is reuse of pretrained models,
checkpointed bounded jobs and rejected-candidate early stopping; we have no
audited end-to-end cost reduction claim. The \textbf{data flywheel} records
failure evidence, accepted repairs and regression checks in reusable modules,
although diagnosis still requires human review. \textbf{Complexity reduction}
comes from explicit geometry/policy contracts and four separable modules:
reconstruction, repair, retargeting and controller training.
\textbf{Innovation} is testing visual and physical validity separately and
retaining simpler controllers when learning regresses. These are demonstrated
research mechanisms, not claims of autonomous production deployment.

Code: \url{https://github.com/isyed/v2d-grounding}; reusable modules:
\url{https://github.com/isyed/real2sim2real}. Exact producing commits, submission
IDs and reproduction status accompany this report. No real-robot evaluation
or generalization beyond the challenge episodes is claimed.
\bibliographystyle{plainnat}
\bibliography{references}
\end{document}
